% concatenated from main.tex, packages.tex, commands.tex, naming.tex
\documentclass[11pt, reqno]{amsart}
\usepackage{amssymb}
\usepackage{amsmath}
\usepackage{mathrsfs}
\usepackage{amsfonts}
\usepackage{color}
%\usepackage{vmargin}
\usepackage{amsthm}
\usepackage{graphicx}
\newtheoremstyle{remark}
  {}{}{}{}{\bfseries}{.}{.5em}{{\thmname{#1 }}{\thmnumber{#2}}{\thmnote{ (#3)}}}
%  \setmarginsrb{20mm}{20mm}{20mm}{20mm}{10mm}{10mm}{10mm}{10mm}
\newtheorem{remark}{Remark}[section]
\newcommand*{\thetheorem}{\Alph{theorem}}

%\RequirePackage{amsthm}
\newcommand{\measurerestr}{%
  \,\raisebox{-.127ex}{\reflectbox{\rotatebox[origin=br]{-90}{$\lnot$}}}\,%
}
 \frenchspacing
\usepackage[utf8]{inputenc} % Required for inputting international characters

%\usepackage[T1]{fontenc} % Output font encoding for international characters

%\usepackage[ansinew]{inputenc}

%\usepackage{mathpazo} % Use the Palatino font by default

%\usepackage[spanish]{babel} % Paquete para escribir en español

\usepackage[width=16cm,height=24cm]{geometry} % Ancho y alto útil de la página.

\usepackage[backend=biber,
%style=chicago-authordate, url=false,doi=false, isbn=false
style=numeric, url=false,doi=false, isbn=false %style=numeric cuando deje de ser un borrador, style=draft para el borrador.
, maxnames=10 %maxnames=10 para que no ponga Autor + et al. si son hasta 10 autores. sorting=nyt para ordenar por autor, año, titulo.
, sorting=nyt]{biblatex}
%\bibliographystyle{abbrv}
\usepackage[autostyle=true]{csquotes} % Required to generate language-dependent quotes in the bibliography


\usepackage{amsfonts,amsmath,amssymb,amsthm}%, scrltx,geometry}
%\geometry{a4paper,twoside,top=4cm,bottom=2cm,left=2cm,right=3cm,headsep=1cm,headheight=3mm}
\usepackage{graphics,graphicx,epsfig}% Required for inserting images

\usepackage{comment}
\usepackage{color}
\usepackage{bbm}
\usepackage{bbold}
\usepackage{enumitem}
\usepackage{appendix}
\usepackage{tikz, tikz-cd}




%\usepackage[dvipsnames,svgnames,table]{xcolor}

\newcommand\boldred[1]{\textcolor{red}{\textbf{#1}}} % Para hacer comentarios visibles. En rojo y en negrita.

\newcommand{\re}{\mathbb{R}}
\newcommand{\nat}{\mathbb{N}}
\newcommand{\be}{\beta}
\newcommand{\la}{\lambda}



\newcommand{\lpv}{{L^{p(\cdot)}(\Omega)}}
\newcommand{\lqv}{{L^{q(\cdot)}(\Omega)}}
\newcommand{\lrv}{{L^{r(\cdot)}(\Omega)}}
\newcommand{\lpvo}{{L^{p(\cdot)}[0,1]}}
\newcommand{\lqvo}{{L^{q(\cdot)}[0,1]}}
\newcommand{\lrvo}{{L^{r(\cdot)}[0,1]}}
\newcommand{\lpvi}{{L^{p(\cdot)}(0,\infty)}}
\newcommand{\lqvi}{{L^{q(\cdot)}(0,\infty)}}
\newcommand{\lrvi}{{L^{r(\cdot)}(0,\infty)}}
\newcommand{\lpvpi}{{L^{p(\cdot)}(0,\pi)}}
\newcommand{\lqvpi}{{L^{q(\cdot)}(0,\pi)}}
\newcommand{\lrvpi}{{L^{r(\cdot)}(0,\pi)}}

\newcommand\Lpv[1]{L^{#1(\cdot)}}

\newcommand{\lpvcomplete}{{L^{p(\cdot)}(\Omega, \Sigma, \mu)}}

\newcommand{\rop}{{\rho_{L^{p(\cdot)}}(}}
\newcommand{\roq}{{\rho_{L^{q(\cdot)}}}}
\newcommand{\ror}{{\rho_{L^{r(\cdot)}}(}}


\newcommand{\essentialrangep}{{R_{p(\cdot)}}}
\newcommand{\essentialrangeq}{{R_{q(\cdot)}}}
\newcommand{\ran}{{R_{p(\cdot)}}}
\newcommand{\esr}{{R_{p(\cdot)}}}

\newcommand{\mi}{{p_{|A_k}^-}}
\newcommand{\ma}{{p_{|A_k}^+}}
\newcommand{\Mii}{{p_{|Q_n}^-}}
\newcommand{\mii}{{p_{|Q}^-}}
\newcommand{\Maa}{{p_{|Q_n}^+}}
\newcommand{\maa}{{p_{|Q}^+}}

\newcommand{\lip}{{Lip(M)}}
\newcommand{\lipdual}{{Lip(M)^*}}
\newcommand{\holder}{{C^{\alpha}(M)}}
\newcommand{\holdervariable}{{C^{\alpha(\cdot)}(M)}}

%\theoremstyle{definition}
\newtheorem{thm}{Theorem}[section]
\newtheorem{defi}[thm]{Definition}
\newtheorem{cor}[thm]{Corollary}
\newtheorem{lemma}[thm]{Lemma}
\newtheorem{prop}[thm]{Proposition}
\newtheorem{exa}[thm]{Example}
\newtheorem{porism}[thm]{Porism}
\newtheorem{conjecture}[thm]{Conjecture}
\newtheorem*{conjecture*}{Conjecture}

\newtheorem{rmk}[thm]{Remark}
\newtheorem{note}[thm]{Note}
%\newtheorem{Rem}[thm]{Remarks}
\newtheorem{obs}[thm]{Observation}


%%%%% Redundancia de nombres completos %%%%%

% \newtheorem{theorem}{Theorem}[section]
% \newtheorem{corollary}[thm]{Corollary}
% \newtheorem{lemma}[thm]{Lemma}
% \newtheorem{proposition}[thm]{Proposition}
% \newtheorem{example}[thm]{Example}
% \newtheorem{porism}[thm]{Porism}

% \newtheorem{remark}[thm]{Remark}
% \newtheorem*{remark*}[thm]{Remark}
% \newtheorem{note}[thm]{Note}
% %\newtheorem{Rem}[thm]{Remarks}
% \newtheorem{observation}[thm]{Observation}
% \newtheorem*{observation*}[thm]{Observation}




%%%%%%%%%% V.2. (Thesis) %%%%%%%%%%


% \newtheorem{definicion}{Definition}[chapter]
% \newtheorem{lema}[definicion]{Lemma}
% \newtheorem{proposicion}[definicion]{Proposition}
% \newtheorem{teorema}[definicion]{Theorem}
% \newtheorem{corolario}[definicion]{Corollary}
% \newtheorem{nota}[definicion]{Remark}
% \newtheorem{observacion}[definicion]{Observation}
% \newtheorem{ejemplo}[definicion]{Example}
% \newtheorem{porisma}[definicion]{Porism}

% \newtheorem*{teorema*}{Theorem}
% \newtheorem*{proposicion*}{Proposition}

%\bibliographystyle{plain}
%\bibliography{references}
\addbibresource{references.bib} % The filename of the bibliography

\title[Banach-Saks property for Hölder spaces]{The weak Banach-Saks property for Hölder spaces}
%\title[my short title]{title}

\author[Górka]{Prezemysław Górka}
\author[Sanchiz]{Mauro Sanchiz$^{*}$}

%\date{October 2025}
\address{Faculty of Mathematics and Information Sciences\\
Warsaw University of Technology\\
Pl. Politechniki 1, 00-661 Warsaw, Poland}
\email{ przemyslaw.gorka@pw.edu.pl}
\address{Departamento de An{\'a}lisis Matem{\'a}tico y Matem\'atica Aplicada, Facultad de Matem{\'a}ticas, Universidad Complutense, 28040 Madrid, Spain}
\email{ msanchiz@ucm.es}
\thanks{$^{*}$ Supported by NAWA under the Ulam postoctoral program 2024/1/00064 and partially supported by the UNED research project: 2025/00145/001 ``Operadores, retículos y estructura de espacios de Banach''}
\keywords {Hölder spaces; Banach-Saks property.}
\subjclass[2020]{
46B26,  	%Nonseparable Banach spaces
%46B40  	%Ordered normed spaces [See also 46A40, 46B42]
%46B42  	%Banach lattices [See also 46A40, 46B40]
%46B45  	%Banach sequence spaces [See also 46A45]
46E30, %Spaces of measurable functions (Lp-spaces,Orlicz spaces, Kothe function spaces, Lorentz ¨spaces, rearrangement invariant spaces, idealspaces, etc.)
 46E15. %Banach spaces of continuous, differentiable or analytic functions
}


\begin{document}

\begin{abstract}
We investigate the weak Banach--Saks property in the setting of H\"older spaces over metric spaces. We show that, for every infinite metric space $(M,d)$ and every $\alpha \in (0,1]$, the H\"older space $C^{\alpha}(M)$ fails to have the weak Banach--Saks property. 


%We show that the Banach space of Lipschitz functions in $[0,1]$ is not weakly-Banach-Saks. %To do so, we construct non trivial functionals $\phi\in Lip(M)^{*}$ (explicitly in $span[f_n]$ for some sequence $(f_n)$ in $Lip(M)$ and non explicitly in $Lip(M)^{*}$ by Hahn Banach's Theorem.
%We prove the some for the Hölder spaces $C^{\alpha}(M)$ and most variable order Hölder spaces $C^{\alpha(\cdot)}(M)$ (when $\alpha^->1$).
\end{abstract}


\maketitle %Crea la cabecera del título



\bigskip
\section{Introduction}
A Banach space $X$ is said to have the \emph{weak Banach--Saks property} if every weakly convergent sequence $x_n \to x$ in $X$ contains a subsequence $(x_{n_k})$ whose Ces\`aro means converge in norm to $x$, that is,
\[
\Big\| \frac{1}{N} \sum_{k=1}^N x_{n_k} - x \Big\|_X \longrightarrow 0 
\quad \text{as } N \to \infty.
\]

In their seminal paper \cite{banach_sur_1930}, Banach and Saks proved that the spaces $L_p[0,1]$ (for $1 < p < \infty$) and $\ell_p$ possess the weak Banach--Saks property. They also presented an incorrect argument claiming that $L_1[0,1]$ fails to have this property. This was later corrected by Szlenk \cite{szlenk_sur_1965}, who showed that $L_1[0,1]$ is in fact weakly Banach--Saks. Furthermore, Schreier \cite{schreier_gegenbeispiel_1930} proved that $C[0,1]$ does not have the weak Banach--Saks property. Finally, Farnum \cite{farnum_banach-saks_1974} established that for a compact metric space $M$, the space $C(M)$ is weakly Banach--Saks if and only if
\[
M^{(\omega)} = \emptyset, 
\quad \text{where } 
M^{(\omega)} = \bigcap_{n=1}^{\infty} M^{(n)},
\]
and $M^{(n)}$ denotes the $n$-th derived set of $M$.

Since then, occasional works on the weak Banach--Saks property have appeared in different context on Banach spaces. More recently, quantifications of the Banach--Saks properties have been studied in \cite{bendova_quantification_2015, silber_quantification_2023}.

The main objective of this paper is to show that, for every infinite metric space $M$, the H\"older space $C^{\alpha}(M)$ does not have the weak Banach--Saks property. This result complements the existing characterizations for spaces of continuous functions and provides a natural extension of earlier work to the H\"older setting.

The paper is organized as follows. In Section~2, we recall several definitions and auxiliary results. We provide a direct proof that $\ell_\infty$ does not have the weak Banach--Saks property. As corollaries, we obtain characterizations of measure spaces $(X,\Sigma,\mu)$ and metric spaces $(M,d)$ for which $L_\infty(X,\Sigma,\mu)$ and $C_b(M)$, respectively, are weakly Banach--Saks. Section~3 contains the proof that $C^{\alpha}(M)$ has the weak Banach--Saks property if and only if the metric space $(M,d)$ is finite.



\section{Preliminaries}
Let $(M,d)$ be a metric space, and let $C_b(M)$ denote the space of all bounded continuous functions on $M$. We equip $C_b(M)$ with the supremum norm
\[
\|f\|_{C(M)} := \sup_{x \in M} |f(x)|.
\]
Fix $0 < \alpha \le 1$. A function $f \in C_b(M)$ is said to be \emph{Hölder continuous of order $\alpha$} if
\[
\rho_{\alpha, M}(f) := \sup_{\substack{x,y \in M \\ x \neq y}} 
\frac{|f(x) - f(y)|}{d(x,y)^\alpha} < \infty.
\]
We define
\[
C^\alpha(M) := \{ f \in C_b(M) : \rho_{\alpha, M}(f) < \infty \},
\]
and equip it with the norm
\[
\|f\|_{C^\alpha (M)} := \|f\|_{C(M)} + \rho_{\alpha, M}(f).
\]
The spaces $C_b(M)$ and $C^\alpha(M)$ are Banach spaces. 

By the definition of the completion of a metric space, we obtain the following.
\begin{prop} \label{iso}
    Let $(M,d)$ be a metric space, and let $(\widehat{M},\widehat{d})$ be its completion. Fix $\alpha \in (0,1]$. Then the following $C^\alpha(\widehat{M}) \cong C^\alpha(M)$ isometric isomorphism holds.
\end{prop}
\begin{proof}
Let $(\widehat{M},\widehat{d})$ denote the completion of $(M,d)$, and let 
$i: M \to \widehat{M}$ be an isometric embedding with dense image. 
Define the operator
\[
I: C^{\alpha}(\widehat{M}) \to C^{\alpha}(M)
\]
by
\[
I(f)(x) = f(i(x)), \qquad \text{for all } f \in C^{\alpha}(\widehat{M}),\ x \in M.
\]

It is immediate that $I$ is linear and injective. Moreover, for every 
$f \in C^{\alpha}(\widehat{M})$,
\[
\|I(f)\|_{C^{\alpha}(M)} \le \|f\|_{C^{\alpha}(\widehat{M})}.
\]

We now prove that $I$ is surjective. Let $g \in C^{\alpha}(M)$ be given and define 
$\widetilde{f}: i(M) \to \mathbb{R}$ by
\[
\widetilde{f}(\hat{x}) = g(i^{-1}(\hat{x})).
\]
For any $\hat{x}, \hat{y} \in i(M)$ we obtain
\[
|\widetilde{f}(\hat{x}) - \widetilde{f}(\hat{y})|
= |g(i^{-1}(\hat{x})) - g(i^{-1}(\hat{y}))|
\le \rho_{\alpha,M}(g)\,d^{\alpha}(i^{-1}(\hat{x}),i^{-1}(\hat{y}))
=\rho_{\alpha,M}(g)\,\widehat{d} ^{\alpha}(\hat{x},\hat{y}),
\]
and
\[
|\widetilde{f}(\hat{x})| \le \|g\|_{C(M)}.
\]
Thus, $\widetilde{f}$ is Hölder continuous on the dense subset 
$i(M) \subset \widehat{M}$ with the same Hölder seminorm and supremum bound as $g$.
Since $i(M)$ is dense in $\widehat{M}$, the function $\widetilde{f}$ admits 
a unique continuous extension $f: \widehat{M} \to \mathbb{R}$. 
Furthermore, this extension satisfies
\[
\rho_{\alpha,\widehat{M}}(f) \le \rho_{\alpha,M}(g),
\qquad
\|f\|_{C(\widehat{M})} \le \|g\|_{C(M)}.
\]
Hence $f \in C^{\alpha}(\widehat{M})$ and, by construction, $I(f)=g$. 
Consequently,
\[
\|f\|_{C^{\alpha}(\widehat{M})}
\le
\|I(f)\|_{C^{\alpha}(M)}.
\]
This completes the proof.
\end{proof}
















Let \(X\) and \(Y\) be Banach spaces with norms \(\|\cdot\|_X\) and \(\|\cdot\|_Y\), respectively. 
We say that \(Y\) embeds into \(X\), denoted \(Y \hookrightarrow X\), if there exists a linear map 
\(T: Y \to X\) and a constant \(C > 0\) such that
\[
C^{-1} \|y\|_Y \leq \|Ty\|_X \leq C \|y\|_Y \quad \text{for all } y \in Y.
\]


\begin{prop}\label{propemb}
Let \(X\) and \(Y\) be Banach spaces. 
If \(X\) is weakly Banach-Saks and \(Y \hookrightarrow X\), 
then \(Y\) is also weakly Banach-Saks.
\end{prop}
\begin{proof}
Let \(T: Y \to X\) be a linear operator. Assume that there exists a constant \(C>0\) such that
\begin{eqnarray}\label{zbi}
C^{-1}\|y\|_{Y} \le \|Ty\|_{X} \le C\|y\|_{Y}
\quad \text{for all } y \in Y.
\end{eqnarray}
Let \((y_n)\subset Y\) be a sequence converging weakly to \(y \in Y\). By linearity and continuity of \(T\), it follows that 
\(T(y_n) \rightarrow T(y)\) weakly in \(X\).
Since \(X\) has the weak Banach--Saks property, there exists a subsequence 
\((y_{n_k})\) such that
\[
\frac{1}{N}\sum_{k=1}^{N} T(y_{n_k}) - T(y)
\longrightarrow 0
\quad \text{in } X.
\]
Using the first inequality in \eqref{zbi}, we deduce that
\[
\frac{1}{N}\sum_{k=1}^{N} y_{n_k} - y
\longrightarrow 0
\quad \text{in } Y.
\]
This completes the proof.

\end{proof}
\begin{thm}\label{zan}
    $\ell_\infty$ is not weakly Banach-Saks.
\end{thm}
\begin{proof}
Although the theorem can be deduced using advanced results, we present a more direct approach for clarity. 
Indeed, since \(C[0,1]\) is separable, Theorem 2.5.7 in \cite{bookA} provides an embedding $
C[0,1] \hookrightarrow \ell_\infty,$
and, by Schreier's Theorem \cite{schreier_gegenbeispiel_1930}, the space \(C[0,1]\) is not weakly Banach-Saks. 
Consequently, Proposition~\ref{propemb} implies that \(\ell_\infty\) is not weakly Banach-Saks. 
While this argument is valid, it relies on relatively heavy machinery; 
to provide a self-contained and more transparent exposition, we shall give the following direct proof following the ideas of Schreier and Farnum.

Let $\mathcal{A} \subset 2^{\mathbb{N}}\setminus \{\emptyset\}$ be the family of maximal Schreier sets \cite{BEANLAND_GOROVOY_HODOR_HOMZA_2024}, defined by
\[
    \mathcal{A}=\{A \subset \mathbb{N}: \vert A\vert=\min(A) \}\setminus\{\emptyset\}.
\] Clearly, $\mathcal{A}$ is a numerable set. Take $T:\mathbb{N}\rightarrow\mathcal{A}$ an enumeration of $\mathcal{A}$, and define the sequence $(u_k)$ in $\ell_\infty$ as follows
$$
u_k(i):=\begin{cases}
    1 \text{ if } k\in T(i)\\
    0 \text{ otherwise.}
\end{cases}
$$
\begin{lemma} \label{jeden}
    The sequence $(u_k)$ is weakly converging in $\ell_\infty$ to $0$.
\end{lemma}
\begin{proof}
We shall need the following criterion for weak convergence in $\ell_{\infty}$.
\begin{prop}[\cite{toland_dual_2020}, Corollary 8.11, rewritten]
    Let $(u_k)$ be a sequence in $\ell_{\infty}$ with $u_k(i) \rightarrow 0$ as $k \rightarrow \infty$ for each $i \in \mathbb{N}$.
    If for every $\alpha>0$ and for all strictly increasing natural sequences $(k_j)$, $(i_n)$ and $(J_n)$, there exists $n\in\mathbb{N}$ and $j\in\{1,2,...,J_n\}$ such that $\vert u_{k_j}(i_n)\vert\leq \alpha$, then $u_k \rightarrow 0$ weakly in $\ell_{\infty}$.
\end{prop}
    To apply the above proposition we fix $\alpha>0$, and strictly increasing sequences of natural numbers $(k_j)$, $(i_n)$ and $(J_n)$. Take $n>k_1$. If $k_1\not\in T(i_n)$, then $u_{k_1}(i_n)=0<\alpha$ and the proof ends, so suppose that $k_1\in T(i_n)$. This forces $\vert T(i_n)\vert=\min(T(i_n))\leq k_1$. Now, since $k_1<n\leq J_n$ and $k_1<k_2<...<k_{J_n}$, there must exist at least some $j\in\{1,...,J_n\}$ such that $k_j\not\in T(i_n)$.
    Hence, $\vert u_{k_j}(i_n)\vert=0<\alpha$. This concludes that $u_k\rightarrow0$ weakly in $\ell_{\infty}$.
\end{proof}

\begin{lemma}\label{dwa}
    The sequence $(u_k)$ has no subsequences whose Cèsaro means converge in $\ell_\infty$ to $0$. 
\end{lemma}

\begin{proof}
        Let us fix any strictly increasing sequence $(k_j)$ of natural numbers,
    observe that $k_{N+1}>N$. Let 
    \[
    A_N:=\{k_{N+1},...,k_{N+N},k_{N+N+1},...,k_{{N+k_{N+1}}}\}.
    \]
    Then the set $\{k_{N+1}, k_{N+2},..., k_{N+N}\}$ is contained in $A_N$. Since $A_N \in \mathcal{A}$, there exists $i_0$ with $T(i_0)=A$, and in consequence $u_{k_j}(i_0)=1$ for every $j\in\{N+1,...,N+N\}$. This leads to
    $$
    \left\| \frac{1}{2 N} \sum_{j=1}^{2N} u_{k_j}\right\|_{\ell_\infty}\geq\frac{1}{2 N} \sum_{j=1}^{2N} u_{k_j}(i_0)\geq \frac{1}{2 N} \sum_{j=N+1}^{2N} u_{k_j}(i_0)=\frac{1}{2}.
    $$
    In conclusion, since for any $N$ we have $\lVert \frac{1}{2 N} \sum_{j=1}^{2N} u_{k_j}\rVert_{\ell_\infty}\geq \frac{1}{2}$, the subsequence $(u_{k_j})$ is not Cèsaro null.
\end{proof}
Now, by Lemma \ref{jeden} and Lemma \ref{dwa} the proof of the theorem follows.
\end{proof}
As corollaries, we get the following two results.
\begin{thm} Let $(X,\Sigma,\mu)$ be measure space, then $L_\infty(X,\Sigma,\mu)$ is not weakly Banach-Saks if and only if there exists a sequence $(A_n)$ of disjoint measurable sets of positive measure.
\end{thm}

\begin{proof}

    $(\Longrightarrow)$ By Proposition~\ref{propemb} and Theorem~\ref{zan}, it is sufficient to prove that $\ell_\infty \hookrightarrow L_\infty(X,\Sigma,\mu).$

Let \((A_n)\) be a sequence of pairwise disjoint measurable sets of positive measure. 
Define an operator \( T : \ell_\infty \to L_\infty(X,\Sigma,\mu) \) by
\[
T\big(a\big) = \sum_{n=1}^{\infty} a(n) \chi_{A_n}, 
\qquad a \in \ell_\infty.
\]
Then \(T\) is a linear isometry. Consequently, \(\ell_\infty\) embeds isometrically 
into \(L_\infty(X,\Sigma,\mu)\), as claimed.
    
       
    $(\Longleftarrow)$ Assume, for the sake of contradiction, that the statement is false. 
Then there exists a natural number \( N \in \mathbb{N} \) and measurable, 
pairwise disjoint sets \( A_1, \dots, A_N \) of positive measure such that $
\mu\left( X \setminus \bigcup_{i=1}^{N} A_i \right) = 0,$
and, moreover, for every \( i \in \{1, \dots, N\} \) and every measurable subset 
\( A \subset A_i \), we have $\mu(A)\,\mu(A_i \setminus A) = 0.$
Then any $f \in L_\infty(X, \Sigma, \mu)$ must be constant almost everywhere on each $A_i$ for $i \in \{1, \dots, N\}$. 
Indeed, if $f \in L_\infty(X,\Sigma,\mu)$ were not constant on $A_i$ for some $i$, then $\operatorname{ess\,sup}_{A_i} f > \operatorname{ess\,inf}_{A_i} f,$
which yields two disjoint measurable subsets of $A_i$ of positive measure — a contradiction. Hence every $f$ is constant on each $A_i$. Thus $L_\infty(X,\Sigma,\mu)$ is finite-dimensional, and therefore it has the weak Banach--Saks property.
\end{proof}
\begin{thm}Let $(M,d)$ be a metric space. Then the following statements hold:
\begin{enumerate}
    \item If $(M,d)$ is compact, the space $C_b(M)$ is weakly Banach--Saks if and only if $M^{(\omega)} = \emptyset$.
    \item If $(M,d)$ is not compact, the space $C_b(M)$ is not weakly Banach--Saks.
\end{enumerate}
\end{thm}
\begin{proof}
Since $(1)$ is just the main result of Farnum \cite{farnum_banach-saks_1974} we shall prove $(2)$. Since $(M,d)$ is not compact, there exists a sequence $(x_n,\varepsilon_n) \subset X \times (0,\infty)$ such that the sequence $(x_n)$ has no convergent subsequence, $\varepsilon_n \to 0$, and
\[
B(x_i,\varepsilon_i) \cap B(x_j,\varepsilon_j) = \emptyset %\varnothing 
\quad \text{for all } i \neq j.
\]
For each $n \in \mathbb{N}$, let $\varphi_n : M \to [0,1]$ be a continuous function satisfying the following conditions\footnote{For example, one may define $\varphi_n(x)=\max\left\{1-\frac{d(x,x_n)}{\varepsilon_n},\,0\right\}$.}
\[
\varphi_n(x_n) = 1 
\quad \text{and} \quad 
\varphi_n \equiv 0 \text{ on } M \setminus B(x_n,\varepsilon_n).
\]
Define the linear operator $T \colon \ell_\infty \to C_b(M)$ by
\[
T\big(a\big) = \sum_{n=1}^{\infty} a(n) \varphi_n, 
\qquad a \in \ell_\infty.
\]
It follows from the properties of the sequence $(x_n,\varepsilon_n)$ that, for every $a \in \ell_\infty$, the map $x \mapsto T(a)(x)$ is continuous. Since $T$ is a linear isometry, it follows from Theorem~\ref{zan} and Proposition~\ref{propemb} that $C_b(M)$ fails to be weakly Banach--Saks.


\end{proof}










\section{Main result}

\begin{thm} Let $(M,d)$ be a metric space and $\alpha \in (0,1]$, then $C^\alpha(M)$ is weakly Banach-Saks if and only if $M$ is finite.
\end{thm}

%\boldred{The case $\alpha=1$, $C^\alpha(M)=Lip(M)$ is already a consequence of \cite{cuth_structure_2016}} {-TINES RAZON QUE ESTO ES UNA CONSEQUENCIA DE \cite{cuth_structure_2016} Y TEOREMA 2.3 Y PROPOSITION 2.2, PERO EN \cite{cuth_structure_2016} TARABAJARON CON $Lip_0 (M)$??-}
\begin{proof}
First, we recall the following lemma from \cite{cuth_structure_2016}\footnote{It is not explicitly stated in that article as we write it, but it is proven in the proof of its Theorem 3.2, which mainly shows the existence of points a sequence of points $(x_n,y_n)$ satisfying the assumptions of its Lemma 3.1, the properties we are taking.}.

\begin{lemma}\label{ciagi}
Let $(M,d)$ be an infinite complete metric space. Then there exist a constant $K>0$ and a sequence of pairs of points $(x_n,y_n) \subset M$ such that:
\begin{enumerate}[label=(\roman*)]
    \item $x_n \neq y_n$ for all $n \in \mathbb{N}$;
    \item $x_m \notin B(y_n, K \, d(x_n,y_n))$ for all $n,m \in \mathbb{N}$ (in particular, this implies $K \leq 1$);
    \item $B(y_n, K \, d(x_n,y_n)) \cap B(y_m, K \, d(x_m,y_m)) = \emptyset$ for all $n \neq m$.
\end{enumerate}
\end{lemma}
Clearly, if $M$ is finite, then $C^{\alpha}(M)$ is finite-dimensional, and thus it is weakly Banach-Saks.
Let $(M, d)$ be an infinite metric space. By Proposition~\ref{iso}, the space $C^\alpha(M)$ is isometrically isomorphic to $C^\alpha(\widehat{M})$, where $(\widehat{M},\widehat{d})$ denotes the completion of the metric space $(M,d)$. Therefore, in view of Proposition~\ref{propemb}, we may assume without loss of generality that $(M,d)$ is complete.



By Proposition \ref{propemb} and Theorem \ref{zan} it is enough to show that $\ell_{\infty} \hookrightarrow C^\alpha(M)$.
Let $K>0$ and let $(x_n,y_n)$ be a sequence of pairs of points in $M$ from Lemma \ref{ciagi}.
   Let us define the sequence of functions $f_n :M \rightarrow \mathbb{R}$ as follows 
    $$
    f_n(x):=\max\left\{ \min\left\{1,d^{\alpha}(x_n,y_n)-\frac{d^\alpha(x,y_n)}{K^\alpha}\right\}, 0 \right\},  \,\, x \in M
    $$
     It is trivial that 
     \[
     \lVert f_n\rVert_{C(M)}\leq 1.
     \]
     Moreover, we have
     \begin{eqnarray} \label{oszacowania}
       \rho_{\alpha, M}(f_n)\leq \frac{1}{K^\alpha}.
     \end{eqnarray}
     Indeed, since the maps $\xi \mapsto \max\{\xi, 0\}$ and $\xi \mapsto \min\{1, \xi\}$ are $1$-Lipschitz, we have 
     
    \begin{align*}
        \rho_{\alpha, M}(f_n)={}    &   \sup_{x\neq y} \frac{\vert f_n(x)-f_n(y)\vert}{d^\alpha(x,y)}\\
        \leq &  \sup_{x\neq y}\frac{\left|d^{\alpha}(x_n,y_n)-\frac{d^\alpha(x,y_n)}{K^\alpha} -\left(d^{\alpha}(x_n,y_n)-\frac{d^\alpha(y,y_n)}{K^\alpha}\right)  \right|}{d^\alpha(x,y)}\\
        = & \sup_{x\neq y} \frac{\left|\frac{d^\alpha(x,y_n)}{K^\alpha} -\frac{d^\alpha(y,y_n)}{K^\alpha}  \right|}{d^\alpha(x,y)}.
    \end{align*}
    Next, we use  the inequality $\vert a^\alpha-b^\alpha\vert\leq \vert a-b\vert^\alpha$ for every $0<\alpha\leq 1$ and $a,b>0$ to compute
    \begin{align*}
        \rho_{\alpha, M}(f_n)\leq{}    &   \sup_{x\neq y}\frac{\vert d(x,y_n)-d(y,y_n)\vert^\alpha}{K^\alpha \cdot d^\alpha(x,y)}\\
    \leq{}  &   \sup_{x\neq y} \frac{\vert d(x,y)\vert^\alpha}{K^\alpha \cdot d^\alpha(x,y)}=\frac{1}{K^\alpha},
    \end{align*}
    and this finishes the proof of (\ref{oszacowania}).


   % -NO NECESITAMOS CONSIDERAR ESTOS CASOS (since the maps $\xi \mapsto \max\{\xi, 0\}$ and $\xi \mapsto \min\{1, \xi\}$ are $1$-Lipschitz,)--
    
    % \boldred{(Aún así, hay que echar las mismas cuentas para el modular, en función de si $x,y\in B_n$ o no. Y además tomando $\min\left\{1,d^{\alpha}(x_n,y_n)-\frac{d^\alpha(x,y_n)}{K^\alpha}\right\}$, que no lo habíamos tenido en cuenta antes. Lo de arriba no es suficiente. Lo escribo de nuevo aquí abajo.)} Let us denote $B_n:=B(y_n,K\cdot d(x_n,y_n))$. From the very definition of $f_n$ we have $supp (f_n):=\{x \in M: f_n (x) \neq 0\} = B_n$.

    % \begin{align*}
    %     \rho_{\alpha}(f_n)={}    &   \sup_{x\neq y} \frac{\vert f_n(x)-f_n(y)\vert}{d^\alpha(x,y)}\\
    %     ={} &   \max\left\{\sup_{x\neq y\in B_n} \frac{\vert f_n(x)-f_n(y)\vert}{d^\alpha(x,y)}, \sup_{x\in B_n, y\not\in B_n} \frac{\vert f_n(x)-0\vert}{d^\alpha(x,y)}, \sup_{x\not\in B_n, y\in B_n} \frac{\vert 0-f_n(y)\vert}{d^\alpha(x,y)}, 0 \right\}.
    % \end{align*}
    % To bound the case $x,y\in B_n$, note that 
    % $$
    % 0=\vert 1-1\vert\leq\vert f_n(x)-f_n(y)\vert\leq \left|d^{\alpha}(x_n,y_n)-\frac{d^\alpha(x,y_n)}{K^\alpha} -\left(d^{\alpha}(x_n,y_n)-\frac{d^\alpha(y,y_n)}{K^\alpha}\right)  \right|,
    % $$
    % so we have
    % \begin{align*}
    % \sup_{x\neq y\in B_n} \frac{\vert f_n(x)-f_n(y)\vert}{d^\alpha(x,y)}\leq{} &   
    % \sup_{x\neq y\in B_n} \frac{\vert \frac{d^\alpha(y,y_n)}{K^\alpha}-\frac{d^\alpha(x,y_n)}{K^\alpha}\vert}{d^\alpha(x,y)}\\
    % \intertext{Next, we use  the inequality $\vert a^\alpha-b^\alpha\vert\leq \vert a-b\vert^\alpha$ for every $0<\alpha\leq 1$ and $a,b>0$ to compute }\\
    % \leq{}  &    \sup_{x\neq y\in B_n} \frac{\vert d(x,y_n)-d(y,y_n)\vert^\alpha}{K^\alpha \cdot d^\alpha(x,y)} \leq \sup_{x\neq y\in B_n} \frac{\vert d(x,y)\vert^\alpha}{K^\alpha \cdot d^\alpha(x,y)}=\frac{1}{K^\alpha}.
    % \end{align*}
    
    % For $x\in B_n$ and $y\not\in B_n$ (analogously $x\not\in B_n$ and $y\in B_n$), note that $\min\{1,d^{\alpha}(x_n,y_n)-\frac{d^\alpha(x,y_n)}{K^\alpha}\}\leq d(x_n,y_n)^\alpha-\frac{d^\alpha(x,y_n)}{K^\alpha}$ and that $d(y,y_n)\geq K\cdot d(x_n,y_n)$, hence
    % \begin{align*}
    % \sup_{x\in B_n, y\not\in B_n} \frac{\vert f_n(x)-0\vert}{d^\alpha(x,y)}
    % \leq{} &   \sup_{x\in B_n, y\not\in B_n} \frac{d^\alpha(x_n,y_n)-\frac{d^\alpha(x,y_n)}{K^\alpha}}{d^\alpha(x,y)}
    % \leq \sup_{x\in B_n, y\not\in B_n} \frac{d^\alpha(y,y_n)-d^\alpha(x,y_n)}{K^{\alpha} \cdot d^\alpha(x,y)}\\
    % \leq{}  &    \sup_{x\in B_n, y\not\in B_n} \frac{\vert d(y,y_n)-d(x,y_n)\vert^{\alpha}}{K^{\alpha} \cdot d^\alpha(x,y)}
    % \leq \sup_{x\in B_n, y\not\in B_n} \frac{d^\alpha(x,y)}{K^{\alpha} \cdot d^\alpha(x,y)}
    % =\frac{1}{K^\alpha}.
    % \end{align*}
    
 
 For each $n \in \mathbb{N}$, let us denote $B_n:=B(y_n,K\cdot d(x_n,y_n))$. From the very definition of $f_n$ we have $supp (f_n):=\{x \in M: f_n (x) \neq 0\} = B_n$.
 Observe that, thanks to property $(iii)$, we have $supp(f_n)\cap supp(f_m)=\emptyset$ for $n\neq m$. Hence, for every $x\in M$, there exists at most one $n\in\mathbb{N}$ such that $f_n(x)\neq 0$. Denote $n(x)$ the unique $n$ with $f_n(x)\neq 0$ (or $n(x)=1$ if there is no such $n$). We are ready now to show that $\ell_{\infty}  \hookrightarrow C^{\alpha}(M)$. For this purpose we define the operator $   
    T \colon \ell_{\infty}  \rightarrow C^{\alpha}(M) $ in the following manner
     \begin{align*}
     T(a)(x):=a(n(x))\cdot f_{n(x)}(x), \,\, a\in \ell_{\infty}, x \in M.
    \end{align*}
     It is clearly lineal, and we will show that 
    \begin{eqnarray}\label{trzy}
    \lVert a\rVert_{\ell_\infty}\leq \lVert T(a)\rVert_{C^\alpha(M)}\leq \left(\frac{2}{K^\alpha}+1\right)\lVert a\rVert_{\ell_\infty}, \,\, \text{for}\,\, a\in \ell_\infty.
    \end{eqnarray}
    Let us fix $a\in\ell_\infty$. 
    On the one hand, fix $x,y\in M$. If $n(x)=n(y)$, then by (\ref{oszacowania}) we have
    \begin{align*}
    \vert T(a)(x)-T(a)(y)\vert =& \vert a(n(x)) (f_{n(x)}(x)-f_{n(x)}(y))\vert \\ \leq & \lVert a\rVert_{\ell_\infty} \cdot \rho_{\alpha, M} (f_{n(x)})d^\alpha(x,y)\leq \lVert a\rVert_{\ell_\infty} \cdot \frac{d^\alpha(x,y)}{K^\alpha}.
    \end{align*}
    Hence, assume $n(x)\neq n(y)$. Thus we have $f_{n(x)}(y)=f_{n(y)}(x)=0$, and consequently, by (\ref{oszacowania}), we obtain 
    \begin{align*}
    \vert T(a)(x)-T(a)(y)\vert 
    ={} &    \vert a(n(x))\cdot f_{n(x)}(x) - a(n(y))\cdot f_{n(y)}(y)\vert \\
    = &\vert a(n(x))\cdot (f_{n(x)}(x)-f_{n(x)}(y)) - a(n(y))\cdot (f_{n(y)}(y)-f_{n(y)}(x))\vert\\
    \leq{}  &   |a(n(x))| \vert f_{n(x)}(x)-f_{n(x)}(y) \vert + |a(n(y))| \vert f_{n(y)}(y)-f_{n(y)}(x)\vert\\
        \leq{}  &   |a(n(x))| \rho_{\alpha, M} (f_{n(x)})d^\alpha(x,y) + |a(n(y))|\rho_{\alpha, M} (f_{n(x)})d^\alpha(x,y) \\
    \leq & \lVert a\rVert_{\ell_\infty} \cdot\frac{2 d^\alpha(x,y)}{K^\alpha}.
    \end{align*}
    Therefore,
    \[
     \rho_{\alpha, M}(T(a)) \leq \frac{2}{K^\alpha} \|a\|_{\ell_\infty},
    \]
    and it is clear that 
    \[
    \lVert T(a)\rVert_{C(M)}\leq\lVert a\rVert_{\ell_\infty} \cdot \sup_n\lVert f_n\rVert_{C(M)}\leq\lVert a\rVert_{\ell_\infty}.
    \]

    On the other hand, for every $k\in\mathbb{N}$, $n(y_k)=k$ and $x_k\not\in \bigcup_{l=1}^\infty B_l$. Thus, for every $l \in \mathbb{N}$ we have  $f_l(x_k)=0$ and
    \begin{eqnarray*}
    \vert T(a)(x_k)- T(a)(y_k) \vert&=&|T(a)(y_k)|=|a(n(y_k))\cdot f_{n(y_k)}(y_k)|\\
    &= &|a(k) \cdot f_k(y_k)|=|a(k)| \cdot \min\{1,d^{\alpha}(x_k,y_k)\}.
    \end{eqnarray*}
    This shows that either $|a(k)|\leq\lVert T(a)\rVert_{C(M)}$ or $|a(k)|\leq \rho_{\alpha, M}(T(a))$. Since $k$ is arbitrary, we get
    \[
    \lVert a\rVert_{\ell_\infty}\leq \rho_{\alpha, M}(T(a)) + \lVert T(a)\rVert_{C(M)} = \lVert T(a)\rVert_{C^\alpha(M)},
    \]
    this finishes the proof of (\ref{trzy}) and the whole proof follows.
\end{proof}








\printbibliography%[heading=bibintoc]




\end{document}
